\section{Zadanie 13}
\autor{Oskar Prusik}
\textbf{Teza:} Niech $f, g : \mathbb{R} \to \mathbb{R}$ będą mierzalne. Czy ich suma również jest mierzalna?

\subsection*{Rozwiązanie}


Niech $\mathcal{G} = \{ (-\infty, c) : c \in \mathbb{R} \}$. 
Wiadomo, że $\sigma(\mathcal{G}) = \mathcal{B}(\mathbb{R})$.
(co łatwo wynika z warunków na $\sigma$-algebrę oraz faktu, że $\sigma(\mathcal{A}_0) = \mathcal{B}(\mathbb{R})$)

Zgodnie z zadaniem 10, aby wykazać mierzalność $f + g$, wystarczy sprawdzić, że dla każdego $G \in \mathcal{G}$:
\[
(f + g)^{-1}(G) \in \mathcal{B}(\mathbb{R}).
\]

Weźmy dowolne $c \in \mathbb{R}$ i rozważmy zbiór:
\[
A = \{ x \in \mathbb{R} : f(x) + g(x) < c \} = (f + g)^{-1}((-\infty, c)).
\]

Zauważmy, że:
\[
f(x) + g(x) < c \quad \iff \quad f(x) < c - g(x).
\]
Istnieje wtedy liczba wymierna $r \in \mathbb{Q}$ taka, że:
\[
f(x) < r < c - g(x).
\]
Zatem:
\[
A = \bigcup_{r \in \mathbb{Q}} \left( \{ x : f(x) < r \} \cap \{ x : g(x) < c - r \} \right).
\]

Ponieważ $f$ i $g$ są mierzalne, zbiory:
\[
\{ x : f(x) < r \} \in \mathcal{B}(\mathbb{R}), \quad \{ x : g(x) < c - r \} \in \mathcal{B}(\mathbb{R}).
\]
Przekrój dwóch zbiorów borelowskich jest borelowski, a przeliczalna suma zbiorów borelowskich jest borelowska. 
Zatem $A \in \mathcal{B}(\mathbb{R})$.

Wnioskujemy, że $(f + g)^{-1}((-\infty, c)) \in \mathcal{B}(\mathbb{R})$ dla każdego $c \in \mathbb{R}$, 
czyli warunek mierzalności jest spełniony na generatorze $\mathcal{G}$. 
Z zadania 10 wynika, że $f + g$ jest mierzalna.


\subsection*{Odpowiedź}
\[
\boxed{\text{Suma też jest mierzalna.}}
\]
